\section*{Chapitre \ref{ch:iv1:statistics} --- Statistiques}
\markboth{\protect\omTocMark{}{Réponses~: Statistiques}}{\protect\omTocMark{}{Réponses~: Statistiques}}

\begin{solution}{iq:iv1:statistics:1}
L'approximation normale donne $z = 2$, une p-valeur d'environ $0.046$, et rejette l'équité. La p-valeur binomiale exacte vaut environ $0.057$, et elle ne rejette pas.

Sous une pièce équilibrée, les faces piles ont une moyenne de 50 et un écart type $\sqrt{100 \cdot \frac14} = 5$\index{approximation normale de la binomiale}, ce qui donne $z = 2$. L'écart vient de ce que le compte est un nombre entier~: la queue exacte commence à 60 et l'inclut tout entier, donc la comparaison équitable utilise $59.5$, ce qui donne $z = 1.9$ et une p-valeur de $0.057$, en accord avec la valeur exacte.

\iqwrong{Traiter $0.05$ comme une vérité tranchée, et affirmer que les deux méthodes se contredisent au sujet de la pièce. Les indices sont limites dans les deux cas, et c'est là le constat honnête.}
\end{solution}

\begin{solution}{iq:iv1:statistics:2}
L'intervalle vaut $0.1~\% \pm 0.098~\%$, soit environ $[0.002~\%, 0.198~\%]$. Diviser la largeur par deux demande quatre fois plus de données~: 1600 jours.

L'erreur type vaut $1\%/\sqrt{400} = 0.05\%$\index{erreur type de la moyenne}, et $1.96$ fois cette valeur donne $0.098\%$. Les erreurs types décroissent en $\sqrt{n}$, donc diviser la largeur par deux demande un échantillon quatre fois plus grand.

\iqwrong{Affirmer qu'il faut deux fois plus de données. La loi en racine carrée dit quatre fois, et c'est pour cela que gagner en précision devient vite coûteux.}
\end{solution}

\begin{solution}{iq:iv1:statistics:3}
Environ 36~\% des idées retenues sont des fausses bonnes idées. Et une p-valeur de $0.05$ ne signifie pas qu'il y a 5~\% de chances que l'hypothèse nulle soit vraie.

Sur 1000 idées\index{fausses découvertes}, 900 sont de fausses bonnes idées, et $0.05$ d'entre elles passent, soit 45. Sur les 100 bonnes idées, 80 passent. La fraction de fausses bonnes idées parmi celles qui passent vaut donc $45/125 = 0.36$.

Une p-valeur est la probabilité d'observer des données aussi extrêmes si l'hypothèse nulle est vraie. La probabilité de l'hypothèse nulle sachant les données a besoin d'une loi a priori, comme le montre ce calcul. Si toutes les idées étaient mauvaises, chaque idée retenue serait une fausse bonne idée, quelle que soit la p-valeur.

\iqwrong{Affirmer qu'une idée retenue donne 95~\% de chances d'être bonne. Cela inverse le conditionnement, et ignore à quel point les bonnes idées sont rares.}
\end{solution}

\begin{solution}{iq:iv1:statistics:4}
Diviser par $n - 1$ corrige le fait que l'on mesure les écarts à la moyenne empirique, et non à la vraie moyenne. La version naïve en $n$ donne en moyenne $\frac{n - 1}{n}\sigma^2$, soit $0.8\sigma^2$ pour des échantillons de taille 5.

Partons de
\[
\sum (X_i - \bar{X})^2 = \sum (X_i - \mu)^2 - n(\bar{X} - \mu)^2.
\]
La moyenne du premier terme vaut $n\sigma^2$, et celle du second vaut $n \cdot \sigma^2/n = \sigma^2$\index{correction de Bessel}. La somme a donc pour moyenne $(n - 1)\sigma^2$. La moyenne empirique se trouve plus près des données que la vraie moyenne, si bien que les écarts qu'on mesure par rapport à elle sont trop petits. Cela ne demande pas la normalité, et une simulation sur des lois exponentielles le confirme.

\iqwrong{Affirmer que l'on perd un point de données. Le terme «~degré de liberté~» est un bon nom, mais il faut savoir montrer le calcul.}
\end{solution}

\begin{solution}{iq:iv1:statistics:5}
Au moins un test passe avec une probabilité $1 - 0.95^{20} \approx 0.64$. Un seuil d'environ $0.0025$ par test maintient la probabilité à environ 5~\%.

Exactement, le seuil par test vaut $1 - 0.95^{1/20} = 0.00256$\index{tests multiples}. La règle simple de Bonferroni, $0.05/20 = 0.0025$, donne une probabilité familiale de $0.049$. Elle tient même lorsque les tests sont dépendants, car elle n'utilise que la borne de l'union.

\iqwrong{Penser que 5~\% reste 5~\%, parce que chaque test a une erreur de 5~\%. La probabilité qu'au moins un soit faux croît vite avec le nombre de tests.}
\end{solution}

\begin{solution}{iq:iv1:statistics:6}
Environ 31 d'entre eux, puisque $1000/2^5 = 31.25$. Sélectionner ceux qui ont un historique sans faute sur cinq ans en trouve 31, et cela ne dit rien du talent.

Chaque gérant a une chance de $1/32$ d'avoir un historique sans faute\index{biais de survie}. Un historique sans faute n'est donc pas rare dans un grand groupe, et l'investisseur verrait une liste de 31 gagnants sans qu'aucun talent n'existe. L'année suivante, chacun a encore 50~\% de chances.

Face à un historique sans faute, la bonne question est combien de personnes ont démarré. Ce nombre ne figure généralement pas sur la page, car les perdants ont fermé.

\iqwrong{Affirmer que $1/32$ vaut environ 3~\%, donc que le résultat est significatif au seuil de 5~\%. Ce test s'applique à un gérant choisi à l'avance, et non au meilleur parmi mille.}
\end{solution}

\begin{solution}{iq:iv1:statistics:7}
La p-valeur bilatérale exacte vaut $0.0414$, par un test des signes.

Si la médiane est nulle, chaque jour est haussier avec une probabilité de $\tfrac12$, quelle que soit la loi\index{test des signes}. Le nombre de jours haussiers suit donc une loi binomiale de paramètres $n = 20$ et $p = \tfrac12$. Alors
\[
P(X \ge 15) = \frac{\binom{20}{15} + \dots + \binom{20}{20}}{2^{20}} = \frac{21\,700}{1\,048\,576},
\]
et le double de cette valeur donne la p-valeur bilatérale. Une simulation avec des tirages de Cauchy donne la même probabilité, comme il se doit. Le test ignore les amplitudes, il perd donc de la puissance sur des données normales, mais il ne se dérègle pas sur les queues épaisses.

\iqwrong{Utiliser un test $t$ sur les rendements. Celui-ci a besoin des amplitudes, qui ne vous ont pas été données, et il serait peu fiable sur 20 points à queues épaisses.}
\end{solution}

\begin{solution}{iq:iv1:statistics:8}
La p-valeur vaut $1/20 = 0.05$. C'est la plus petite valeur que puisse donner un test de permutation avec deux groupes de 3.

Si les étiquettes n'ont aucun effet, chaque choix des 3 valeurs parmi 6 appelées A est aussi probable\index{test de permutation}. Il y a $\binom{6}{3} = 20$ choix possibles. Le choix observé, $\{5, 7, 9\}$, a la plus grande somme possible, 21, et c'est le seul à l'atteindre. La p-valeur vaut donc $1/20$.

C'est le plancher pour deux groupes de 3, puisqu'une seule affectation est la plus extrême possible. À cette taille d'échantillon, on ne peut jamais descendre sous 5~\%, quelle que soit la taille réelle de l'effet.

\iqwrong{Lancer un test $t$, qui donne une valeur bien plus petite. Il tire cette valeur d'une hypothèse de normalité des valeurs, que l'on ne peut pas vérifier avec 6 points.}
\end{solution}

\begin{solution}{iq:iv1:statistics:9}
L'intervalle normal vaut $1 \pm 0.877$, et l'intervalle de $t$ vaut $1 \pm 1.241$. La différence vient de ce que l'écart type est estimé, et elle cesse de compter vers 30 points, où la valeur de $t$ vaut $2.05$.

L'erreur type vaut $1/\sqrt{5} = 0.447$\index{loi de Student}. En multipliant par $1.96$ et par $2.776$, on obtient les deux intervalles. Lorsque $\sigma$ est estimé à partir de 5 points, une estimation basse est fréquente, ce qui rend le rapport grand plus souvent qu'avec une variable normale. Les queues sont donc plus épaisses. Une simulation montre que l'intervalle normal ne couvre qu'environ $87.8\%$ du temps, et l'intervalle de $t$ $95\%$.

\iqwrong{Affirmer que l'écart vient de données non normales. Ici, les données sont normales. La largeur supplémentaire vient seulement de ce que $\sigma$ n'est pas connu.}
\end{solution}

\begin{solution}{iq:iv1:statistics:10}
La règle empirique donne $1/\sqrt{n} \approx 0.099$ au voisinage de zéro. L'intervalle à 95~\% pour $r = 0.5$ vaut environ $[0.34, 0.63]$.

La transformation de Fisher\index{transformation z de Fisher}, $z = \operatorname{arctanh}(r)$, rend la loi d'échantillonnage presque normale, avec une erreur type de $1/\sqrt{n - 3} = 0.1$, quelle que soit la vraie corrélation. On a donc $z = 0.549 \pm 0.196$. On revient à l'échelle initiale avec $\tanh$, ce qui donne l'intervalle. Une simulation avec une vraie corrélation de $0.5$ confirme un écart type de $0.1$ sur l'échelle transformée, et une couverture de 95~\%.

L'intervalle n'est pas symétrique autour de $0.5$, car la corrélation est bornée par 1, et la transformation dilate les valeurs proches de $\pm 1$. Une estimation de $0.5$ dispose de plus de marge en dessous qu'au-dessus.

\iqwrong{Utiliser $0.5 \pm 2 \cdot 0.099$. Au voisinage de zéro, cela convient, mais en $0.5$ l'erreur type réelle vaut environ $\frac{1 - r^2}{\sqrt{n}} \approx 0.074$, bien plus petite, et l'intervalle est asymétrique.}
\end{solution}

\begin{solution}{iq:iv1:statistics:11}
La borne exacte vaut $1 - 0.05^{1/100} \approx 0.0295$, et la règle empirique donne $3/n = 0.03$. L'estimation naïve vaut 0, ce qui affirme qu'une panne est impossible.

Trouver la plus grande probabilité de panne $p$ qui reste compatible avec l'absence de pannes\index{règle de trois}. Il faut $(1 - p)^{100} \ge 0.05$. En prenant les logarithmes, on obtient $100p \approx -\ln 0.05 \approx 3$. Donc 100 jours sans incident indiquent que la probabilité est inférieure à environ 3~\% par jour. Cela correspond à une panne environ tous les 33 jours, ce qui n'a rien de rassurant.

\iqwrong{Donner l'estimation ponctuelle 0, ou $1/100$. Il faut un intervalle, et sans panne c'est la borne supérieure qui compte.}
\end{solution}

\begin{solution}{iq:iv1:statistics:12}
Il vous faut 16 ans à un ratio de Sharpe de $0.5$, 4 ans à 1 et 1 an à 2.

Mesurée sur $T$ années, un ratio de Sharpe annualisé a une erreur type voisine de $1/\sqrt{T}$\index{inférence sur le ratio de Sharpe}. C'est la même erreur type que celle d'une moyenne exprimée en unités de son propre écart type, et elle ne dépend pas de la fréquence d'échantillonnage. Le ratio vaut donc $\mathrm{SR}\sqrt{T}$, et il atteint 2 lorsque $T = 4/\mathrm{SR}^2$. Une simulation sur 16 ans de rendements quotidiens confirme une erreur type de $0.25$.

Un historique de trois ans ne permet donc pas de distinguer un Sharpe de $0.5$ de zéro, et aucun historique réaliste ne le permet pour une stratégie lente. Seules les stratégies à haute fréquence, dont le ratio de Sharpe est très élevé, révèlent rapidement leur talent.

\iqwrong{Croire que les données quotidiennes aident, parce qu'elles sont très nombreuses. Un ratio de Sharpe est une moyenne divisée par une volatilité. Seule la volatilité profite de l'échantillonnage plus fin.}
\end{solution}

\begin{solution}{iq:iv1:statistics:13}
Il vous faut $n = (1.96 + 0.84)^2(\sigma/\delta)^2 = 3136$ jours, soit environ $12.4$ ans. Un effet deux fois plus grand n'en demande que le quart~: 784 jours, soit environ $3.1$ ans.

Le test rejette lorsque la moyenne empirique s'écarte de plus de $1.96$ erreurs types de zéro\index{calcul de puissance}. Pour avoir une puissance de 80~\%, l'effet réel doit se trouver à $1.96 + 0.84$ erreurs types de zéro, car $0.84$ est le quantile à 80~\% de la loi normale. On a donc $\delta\sqrt{n}/\sigma = 2.8$, et en élevant au carré, $n = 7.84 \cdot 400$. Une simulation du test confirme une puissance de 80~\%.

\iqwrong{Ne retenir que le $1.96$, ce qui donne environ 1537 jours. À cette taille, on n'a qu'une chance sur deux de détecter l'effet, car la moitié du temps l'estimation tombe en dessous de la vraie valeur.}
\iqwhy{Il faut planifier l'échantillon avant de tester. Un avantage quotidien de 5 points de base avec une volatilité de 1~\% correspond à un Sharpe proche de 0.8, et il faut plus d'une décennie pour le confirmer par cette règle.}
\end{solution}

\begin{solution}{iq:iv1:statistics:14}
Bonferroni rejette un test, et Benjamini--Hochberg en rejette trois. Bonferroni contrôle la probabilité d'au moins un faux rejet, tandis que Benjamini--Hochberg contrôle l'espérance de la proportion de faux rejets parmi les rejets.

Bonferroni compare chaque p-valeur à $0.05/10 = 0.005$. Seule $0.001$ passe ce seuil.

Benjamini--Hochberg\index{procédure de Benjamini--Hochberg} les trie, puis cherche le plus grand $k$ tel que $p_{(k)} \le k \cdot 0.05/10$. Ici $0.001 \le 0.005$, $0.008 \le 0.010$, $0.012 \le 0.015$, mais $0.041 > 0.020$, et aucune valeur suivante ne passe. Les trois premières sont donc rejetées. Si l'on cherche des signaux qu'on validera ensuite, un taux de fausses découvertes est le choix naturel. Si un faux résultat est coûteux, il faut retenir la règle de contrôle du risque global.

\iqwrong{S'arrêter au premier échec d'une règle de montée sans regarder plus loin. Benjamini--Hochberg retient le dernier $k$ qui passe, même si un rang antérieur échoue.}
\end{solution}

\begin{solution}{iq:iv1:statistics:15}
La méthode des moments donne $\hat\theta = 8/4 = 2$ et $\hat{k} = 4/2 = 2$. Le maximum de vraisemblance est plus efficace, surtout pour le paramètre de forme, au prix d'une résolution numérique.

Il faut égaler les moments\index{méthode des moments}. La variance divisée par la moyenne donne $\theta$, puis $k$ vaut la moyenne divisée par $\theta$.

La méthode est simple, et elle fournit une bonne valeur de départ. Mais elle utilise la variance empirique, qui est bruitée pour des données dissymétriques. Sur 300 échantillons de taille 50, la forme obtenue par maximum de vraisemblance présente une variance plus faible que celle obtenue par les moments. Le maximum de vraisemblance est asymptotiquement efficace\index{maximum de vraisemblance}, à condition que le modèle soit juste. On part donc souvent des valeurs des moments, puis on itère.

\iqwrong{Affirmer que les deux méthodes sont équivalentes, puisqu'elles utilisent les mêmes données. Elles ne les exploitent pas de la même façon, et elles ne coïncident que pour des familles particulières, comme la loi de Poisson ou la loi normale.}
\end{solution}

\begin{solution}{iq:iv1:statistics:16}
L'estimateur du maximum de vraisemblance est $\hat\lambda = 1/\bar{X}$. Il est biaisé à la hausse~: $E[\hat\lambda] = \frac{n}{n - 1}\lambda$. La version sans biais vaut $\frac{n - 1}{n\bar{X}}$, soit $0.4$ contre $0.5$ pour $n = 5$ et $\bar{X} = 2$.

La log-vraisemblance vaut $n\ln\lambda - \lambda\sum x_i$\index{maximum de vraisemblance}. En annulant sa dérivée, on obtient $\lambda = n/\sum x_i$.

La somme suit une loi Gamma de paramètre de forme $n$, donc $E[1/\sum x_i] = \lambda/(n - 1)$, ce qui donne le biais. En clair~: $1/x$ est convexe, si bien que, par l'inégalité de Jensen\index{inégalité de Jensen}, la moyenne de $1/\bar{X}$ dépasse $1/E[\bar{X}]$. Une simulation de paramètre 3 avec $n = 5$ donne $3.75$ pour le premier et 3 pour le second.

\iqwrong{Croire qu'une moyenne sans biais donne un taux sans biais. Une fonction non linéaire d'un estimateur sans biais est biaisée.}
\end{solution}

\begin{solution}{iq:iv1:statistics:17}
La borne de Cramér--Rao vaut $\lambda/n$, et $\mathrm{Var}(\bar{X}) = \lambda/n$, donc la borne est atteinte. Pour $\lambda = 2$, les variances asymptotiques de la probabilité de zéro valent $0.0366/n$ pour $e^{-\bar{X}}$ et $0.117/n$ pour la proportion de jours nuls. Ainsi, $e^{-\bar{X}}$ est environ $3.2$ fois meilleur, si les données suivent bien une loi de Poisson.

L'information de Fisher par jour vaut $1/\lambda$\index{borne de Cramér--Rao}, puisque le score vaut $x/\lambda - 1$ et que sa variance vaut $\lambda/\lambda^2$. La borne est donc $\lambda/n$.

Pour $e^{-\bar{X}}$, il faut utiliser la méthode delta~: la pente vaut $-e^{-\lambda}$, ce qui donne $\lambda e^{-2\lambda}/n$. La proportion de zéros est une proportion binomiale, de variance $e^{-\lambda}(1 - e^{-\lambda})/n$. Des simulations sur 200 jours confirment les deux résultats.

La première approche repose sur le modèle. Avec des comptages surdispersés de moyenne 2, et une vraie probabilité de zéro égale à $1/3$, elle annoncerait $e^{-2} \approx 0.135$, très loin de la vérité, alors que la proportion reste correcte. L'efficacité a donc un prix, qui est l'hypothèse de modèle.

\iqwrong{Préférer la proportion parce qu'elle est «~non paramétrique~», et la juger meilleure sans nuance, ou l'inverse. Chacune est le bon choix selon le degré de confiance accordé au modèle.}
\end{solution}

\begin{solution}{iq:iv1:statistics:18}
Environ $1 - 1/e \approx 63.2~\%$. Le bootstrap estime la variabilité d'échantillonnage d'une statistique, en traitant l'échantillon comme s'il était la population.

Un point donné est absent d'un tirage avec probabilité $1 - 1/n$\index{bootstrap}, et de tous les $n$ tirages avec probabilité $(1 - 1/n)^n \to 1/e$. Donc $63.2\%$ des points apparaissent, certains plusieurs fois.

L'idée~: on aimerait tirer de nombreux échantillons dans la vraie distribution et observer comment la statistique varie. Ce n'est pas possible, alors on tire dans l'échantillon lui-même, qui est la meilleure estimation dont on dispose de cette distribution. Pour une moyenne, l'erreur type bootstrap coïncide avec la formule $s/\sqrt{n}$, ce qu'une simulation confirme. Pour une médiane ou un ratio de Sharpe, où il n'existe parfois pas de formule simple, il en fournit une sans aucun calcul formel.

\iqwrong{Affirmer que le bootstrap crée des données supplémentaires. Il ne crée aucune information sur la population. Il montre seulement de combien la statistique varie par rééchantillonnage.}
\end{solution}

\enlargethispage{-1\baselineskip}%
\begin{solution}{iq:iv1:statistics:19}
La variance est trois fois plus grande~: $(1 + \rho)/(1 - \rho) = 3$, donc l'erreur type est $\sqrt{3} \approx 1.73$ fois plus grande. Un bootstrap classique donne le chiffre naïf. Il faut utiliser un bootstrap par blocs.

La variance d'une moyenne vaut\index{autocorrélation et variance de la moyenne}
\[
\frac{\sigma^2}{n}\Bigl(1 + 2\sum_{k \ge 1} \rho^k\Bigr) = \frac{\sigma^2}{n} \cdot \frac{1 + \rho}{1 - \rho}.
\]
Avec $\rho = 0.5$, la somme vaut $1 + 2 \cdot 1 = 3$. Des trajectoires simulées de longueur 2000 montrent ce facteur.

Un bootstrap classique tire les points indépendamment, il détruit donc l'autocorrélation et annonce $\sigma^2/n$, comme le confirme une simulation sur une trajectoire. Un bootstrap par blocs\index{bootstrap par blocs} tire des suites de points voisins, plus longues que la mémoire de la série, ici des blocs de 50, et cela retrouve l'essentiel du facteur.

\iqwrong{Utiliser le bootstrap classique comme remède universel. Il suppose des tirages indépendants, et le résultat est trop étroit sur toute donnée autocorrélée.}
\iqfollow{Et pour une autocorrélation négative~?}{Le facteur est inférieur à 1. Pour $\rho = -0.5$, il vaut $1/3$, de sorte que la moyenne est connue avec plus de précision que ne le suggère la formule naïve.}
\end{solution}

\begin{solution}{iq:iv1:statistics:20}
Sous normalité, les erreurs types valent $\sqrt{6/n} \approx 0.077$ pour l'asymétrie et $\sqrt{24/n} \approx 0.155$ pour le kurtosis en excès. Jarque--Bera vaut $\frac{1000}{6}\bigl(0.25 + \frac94\bigr) \approx 417$, à comparer à un chi-deux à 2 degrés de liberté, dont le quantile à 95~\% vaut $5.99$.

La statistique est la somme des deux $z$-scores élevés au carré\index{test de Jarque--Bera}, soit $\frac{n}{6}S^2 + \frac{n}{24}K^2$. Des simulations d'échantillons normaux de taille 1000 confirment les deux erreurs types.

Le test rejette facilement sur des données financières, et c'est là sa faiblesse. Avec beaucoup de points, il détecte n'importe quel écart, même minime. Et il indique seulement que les données ne sont pas normales, sans dire pourquoi. Une volatilité stochastique à chocs normaux échoue également au test, dans la simulation plus de 99~\% du temps, donc un rejet n'implique pas des chocs à queues épaisses. Les estimations de l'asymétrie et du kurtosis sont aussi très sensibles à quelques journées extrêmes, il faut donc les examiner avant de s'y fier.

\iqwrong{Prendre le rejet pour une découverte utile. Pour les rendements quotidiens, le test rejette presque toujours, et un graphique des queues en apprend davantage.}
\end{solution}

\begin{solution}{iq:iv1:statistics:21}
Si la loi est normale, la probabilité vaut environ $2.87 \times 10^{-7}$ par jour, soit une fois tous les $13\,800$ ans environ. Si la loi est de Student $t$ à 3 degrés de liberté, la probabilité vaut environ $0.0016$ par jour, soit une fois tous les $2.5$ ans environ. Un tel jour n'a rien d'un miracle. Il montre que le modèle est faux.

Pour la loi normale, $\Phi(-5) \approx 2.87 \times 10^{-7}$, et avec 252 jours de séance par an, on obtient le premier chiffre.

Pour la loi de Student $t$, la variance vaut $\nu/(\nu - 2) = 3$\index{queues épaisses}, donc un mouvement recentré de $-5$ correspond à une valeur brute de $t$ égale à $-5\sqrt{3} = -8.66$. Sa queue vaut environ $0.0016$. Une simulation confirme cette queue.

Il ne faut pas en conclure qu'un événement très rare s'est produit. Il faut comparer la vraisemblance sous chaque modèle. Si les journées à 5 sigmas surviennent tous les quelques ans, c'est la loi normale qu'il faut abandonner, et une limite de stress construite sur elle doit être refaite.

\iqwrong{Affirmer que l'événement est impossible, ou qu'il s'agit simplement de malchance. Une probabilité aussi faible est une affirmation sur un modèle, et la preuve que ce modèle est faux est bien plus probable que l'événement lui-même.}
\end{solution}

\begin{solution}{iq:iv1:statistics:22}
La statistique du chi-deux vaut 8, en dessous de $11.07$, donc l'équilibre n'est pas rejeté (p-valeur d'environ $0.156$). On pourrait tester la seule face 1, avec $z \approx 2.45$ (p-valeur d'environ $0.014$), mais on l'a choisie après avoir regardé les résultats. Il faut donc corriger pour les 6 faces, ce qui donne environ $0.09$.

Chaque effectif attendu vaut 20\index{test du chi-deux}. La statistique vaut
\[
\sum \frac{(O - 20)^2}{20} = \frac{100 + 4 + 0 + 4 + 16 + 36}{20} = 8.
\]
Pour une seule face, l'effectif suit une loi binomiale, d'écart type $\sqrt{120 \cdot \frac16 \cdot \frac56} = 4.08$, donc $z = 10/4.08 = 2.45$.

\iqwrong{Tester la face la plus élevée et rapporter $0.014$. C'est la plus élevée parmi six par construction, donc cette p-valeur est six fois trop belle.}
\end{solution}

\begin{solution}{iq:iv1:statistics:23}
La moyenne vaut $50.5$ et l'écart type environ $4.97$. Trente séries sont à $4.1$ écarts types en dessous de la moyenne, ce qui suggère du momentum ou du partitionnement (clustering). Soixante-cinq sont à $2.9$ au-dessus, ce qui suggère un retournement.

Le nombre de séries est égal à 1 plus le nombre de changements\index{test des séries}. Chacune des 99 paires voisines est un changement avec une probabilité de $\tfrac12$, et ces événements sont indépendants pour une suite équilibrée. Le nombre de changements suit donc une loi binomiale de paramètres $n = 99$, de moyenne $49.5$ et de variance $99/4$. Une simulation confirme ces deux valeurs.

Il faut noter que le test porte sur l'ordre, et non sur la proportion de jours haussiers. Un marché peut avoir un rendement moyen positif et sembler pourtant parfaitement aléatoire selon ce test.

\iqwrong{Penser qu'une longue série est suspecte en soi. Les longues séries sont courantes dans des données équilibrées, et on a tendance à les sous-estimer. Ce qui compte, c'est le nombre total de séries, comparé à cette moyenne et à cet écart type.}
\end{solution}

\begin{solution}{iq:iv1:statistics:24}
Avec des échantillons indépendants, $z \approx 0.71$, et avec des paires, $z \approx 2.24$. Le couplage par paires est le bon, car les deux algorithmes ont exécuté les mêmes ordres.

Avec des échantillons indépendants, l'erreur type de la différence vaut $\sqrt{(100 + 100)/100} = 1.41$. Avec des paires, la variance de la différence sur chaque ordre vaut $2 \cdot 100 \cdot (1 - 0.9) = 20$\index{test apparié}, de sorte que l'erreur type vaut $\sqrt{20/100} = 0.447$. Elle est donc environ $\sqrt{10} \approx 3.16$ fois plus petite. Une simulation de paires corrélées donne le même chiffre.

La difficulté des ordres est commune aux deux méthodes, elle s'annule donc dans la différence. Pour comparer des méthodes, il faut concevoir le test de façon que le bruit commun puisse s'annuler~: exécuter les deux sur les mêmes ordres, ou sur des ordres voisins affectés au hasard. Ici, les mêmes données passent de «~aucune preuve~» à «~significatif~» par le seul choix du bon protocole.

\iqwrong{Utiliser le test indépendant parce que c'est celui par défaut, et conclure que les deux algorithmes ne peuvent pas être distingués. La formule doit refléter la façon dont les données ont été produites.}
\end{solution}

\begin{solution}{iq:iv1:statistics:25}
Oui, c'est significatif. L'écart vaut 3 avec une erreur type de $\sqrt{2}$, donc $z \approx 2.12$ et la p-valeur vaut environ $0.034$. Avec des erreurs égales, les intervalles se chevauchent jusqu'à ce que l'écart atteigne $3.92$ erreurs types, alors que la significativité n'en exige que $2.77$.

Les deux intervalles sont $0 \pm 1.96$ et $3 \pm 1.96$, qui se chevauchent entre $1.04$ et $1.96$\index{chevauchement des intervalles de confiance}. Mais l'erreur type d'une différence n'est pas la somme des deux erreurs types~: c'est $\sqrt{1 + 1}$. Le non-chevauchement est donc un critère très strict, dont le niveau réel est d'environ $0.006$. Un chevauchement ne signifie pas «~aucune différence~».

De façon générale, avec des erreurs types $s_1$ et $s_2$, les intervalles se chevauchent tant que l'écart est inférieur à $1.96(s_1 + s_2)$, et la différence est significative dès que l'écart dépasse $1.96\sqrt{s_1^2 + s_2^2}$. Une différence significative avec des intervalles qui se chevauchent est donc toujours possible, et les deux seuils diffèrent d'un facteur allant jusqu'à $\sqrt{2}$, atteint lorsque $s_1 = s_2$.

\iqwrong{Affirmer que des intervalles qui se chevauchent signifient l'absence de différence significative. Le non-chevauchement est la règle la plus stricte, d'un facteur allant jusqu'à $\sqrt{2}$ sur l'écart.}
\end{solution}

\begin{solution}{iq:iv1:statistics:26}
La corrélation de Pearson vaut environ $0.59$ et celle de Spearman vaut 1. Spearman donne le résumé le plus juste de cette relation, qui est parfaitement croissante.

Spearman est la corrélation des rangs\index{corrélation de rang}. Les rangs de $y$ sont ceux de $x$, donc elle vaut exactement 1. Pearson mesure une relation linéaire, si bien qu'un seul point loin de la droite la tire vers le bas. Sans ce point, elle vaut elle aussi 1.

Utiliser Spearman lorsque l'ordre compte, ou lorsque les données comportent des valeurs aberrantes ou des queues épaisses. Elle ignore les amplitudes, si bien qu'elle peut manquer le fait que le point le plus élevé compte vraiment. Si l'on dimensionne une position à partir d'un signal, la relation linéaire effective importe, et Pearson sur des données winsorisées peut être plus pertinent. Présenter les deux et, si elles diffèrent, chercher la raison.

\iqwrong{Affirmer que $0.59$ est la vraie force de la relation, ou retirer le point parce qu'il semble étrange. Commencer par chercher pourquoi il est là.}
\end{solution}

\begin{solution}{iq:iv1:statistics:27}
Le meilleur diviseur est $c = n + 1$. Pour $n = 10$, les erreurs quadratiques moyennes valent $0.2222$ pour $c = 9$, $0.19$ pour $c = 10$ et $0.1818$ pour $c = 11$.

Posons $S = \sum(X_i - \bar{X})^2$, de sorte que $S/\sigma^2$ suit une loi du khi-deux à $m = n - 1$ degrés de liberté\index{erreur quadratique moyenne}. Donc $E[S] = m\sigma^2$ et $E[S^2] = (2m + m^2)\sigma^4$. Alors, en unités de $\sigma^4$,
\[
E\Bigl[\Bigl(\frac{S}{c} - 1\Bigr)^2\Bigr] = \frac{2m + m^2}{c^2} - \frac{2m}{c} + 1.
\]
En minimisant par rapport à $1/c$, on obtient $c = \frac{2m + m^2}{m} = m + 2 = n + 1$. La valeur minimale vaut $2/(n + 1)$.

Le meilleur diviseur donne un estimateur biaisé~: ici, il vaut en moyenne $\frac{9}{11}\sigma^2$. C'est un bon compromis, car les estimations de variance ont elles-mêmes une forte variance, et la réduire diminue celle-ci\index{compromis biais-variance}. L'absence de biais est une propriété, qui n'est pas équivalente à la précision. Le résultat dépend de l'aplatissement de la loi normale. Avec une loi de Laplace, une simulation donne un meilleur diviseur supérieur à 12, il faut donc réduire encore davantage sur les queues épaisses.

\iqwrong{Affirmer que $n - 1$ doit être le meilleur, puisqu'il est sans biais. Sans biais signifie juste en moyenne. On peut avoir une erreur quadratique moyenne plus faible en acceptant un petit biais, même si l'estimateur sans biais est le plus proche sur la plupart des échantillons~: ici, $S/11$ ne fait mieux que $S/9$ que dans environ 36~\% des cas.}
\end{solution}

\begin{solution}{iq:iv1:statistics:28}
Rejeter lorsque $\bar{X} > \tfrac12$, ce qui correspond au test du rapport de vraisemblance. Il faut $n = 11$, car $(2 \cdot 1.645)^2 = 10.8$.

D'après le lemme de Neyman--Pearson\index{lemme de Neyman--Pearson}, le test le plus puissant à un niveau quelconque utilise le rapport de vraisemblance. Pour deux lois normales de même variance,
\[
\ln\frac{L_1}{L_0} = \sum_i \Bigl(x_i - \frac12\Bigr) = n\Bigl(\bar{x} - \frac12\Bigr),
\]
si bien que le test ne dépend que de la moyenne empirique. Fixer le seuil en $\tfrac12$ donne des taux d'erreur égaux, par symétrie.

Chaque erreur vaut $P\bigl(Z > \tfrac12\sqrt{n}\bigr)$, car le seuil se situe à une demi-distance entre les moyennes, en unités de $1/\sqrt{n}$. Pour que chaque erreur reste au plus égale à 5~\%, il faut $\tfrac12\sqrt{n} \ge 1.645$. Pour $n = 11$, l'erreur vaut $0.049$, et pour $n = 10$, elle dépasse 5~\%. Une simulation à $n = 11$ confirme ce résultat pour les deux erreurs.

Pour un écart de $\delta$ écarts types, $n \approx 10.8/\delta^2$. Diviser l'écart par deux exige quatre fois plus de données.

\iqwrong{Utiliser un test unilatéral contre la première hypothèse au niveau de 5~\%, avec un seuil de $1.645/\sqrt{n}$. Cela ne fixe qu'un seul taux d'erreur, et ne dit rien de l'autre. Une décision équitable entre deux hypothèses exactes maîtrise les deux.}
\end{solution}

\begin{solution}{iq:iv1:statistics:29}
Le vrai taux de faux positifs est d'environ 19~\%. Sans limite sur le nombre de regards, il grimpe vers 100~\%. Un seuil d'environ $2.56$ par regard ramène ce taux à 5~\%, et c'est l'idée derrière les tests séquentiels par groupes, avec les frontières de Pocock.

Chaque regard est corrélé aux précédents, si bien que le risque ne s'additionne pas complètement, mais il s'accumule tout de même\index{arrêt optionnel dans les expériences}. En simulant les scores $z$ d'une somme cumulée sur 10 regards, on obtient $19\%$. Avec 5 regards, on obtient $14\%$.

Sans limite, la loi du logarithme itéré affirme que le score $z$ d'une marche aléatoire finit par dépasser tout niveau fixe\index{loi du logarithme itéré}. On est donc certain de «~trouver~» un effet si l'on est prêt à attendre assez longtemps. Une simulation à 2000 regards dépasse déjà 40~\%.

Un seuil de $2.56$ par regard, qui correspond à une p-valeur nominale d'environ $0.01$, maintient le total à 5~\% pour 10 regards. D'autres schémas répartissent l'erreur de façon inégale, de sorte que les premiers regards exigent des preuves très solides.

\iqwrong{Penser que regarder en cours de route est sans danger si l'on ne modifie pas la procédure. La décision d'arrêter fait partie de la procédure, et une p-valeur calculée comme si la taille de l'échantillon était fixée est fausse.}
\iqwhy{Dans la recherche comme dans le backtest, arrêter lorsque les résultats semblent bons est le même piège que tester de nombreux signaux. Fixer la taille de l'échantillon avant de commencer, ou utiliser une frontière conçue pour l'arrêt.}
\end{solution}

\begin{solution}{iq:iv1:statistics:30}
Le maximum du bootstrap égale le maximum de l'échantillon avec la probabilité $1 - (1 - 1/n)^n \to 1 - 1/e \approx 0.632$. La vraie loi d'échantillonnage n'a pas une telle masse. Il faut utiliser la loi exacte~: $P(\max/\theta \le u) = u^n$.

Un échantillon bootstrap contient le point le plus grand de l'échantillon, sauf si les $n$ tirages le manquent tous\index{echec du bootstrap@échec du bootstrap}, de sorte que la loi bootstrap possède un atome de $63.2\%$ au maximum de l'échantillon. Tout intervalle bootstrap est donc dégénéré d'un côté. Il ne peut pas refléter le fait que le vrai $\theta$ se situe au-dessus de chaque point observé.

La vraie loi est continue. $n(\theta - \max)/\theta$ tend vers une loi exponentielle de moyenne 1, et des simulations en confirment la moyenne et une valeur de queue. Comme $\max/\theta$ suit la loi exacte ci-dessus, un intervalle unilatéral à 95~\% est $\theta \le \max/0.05^{1/n}$, et pour $n = 1000$, le facteur vaut $1/0.9970$.

Pourquoi cela échoue~: le bootstrap fonctionne lorsque la statistique est une fonction lisse de la distribution empirique. Une valeur extrême dépend de ce qui se passe au bord, et l'échantillon ne contient aucune information au-delà de ce bord. Un sous-échantillon de taille $m \ll n$ fonctionne mieux, tout comme un bootstrap paramétrique.

\iqwrong{Affirmer que le bootstrap fonctionne toujours si l'on prend assez de rééchantillons. Davantage de rééchantillons réduisent seulement l'erreur de Monte-Carlo. Ils ne réparent pas une méthode inadaptée à cette statistique.}
\iqwhy{Les mesures de risque de queue sont proches de ce cas. Faire un bootstrap sur la pire journée ou sur le quantile extrême d'une série de rendements sous-estime à quel point les choses peuvent mal tourner, pour la même raison.}
\end{solution}

\begin{solution}{iq:iv1:statistics:31}
Sous normalité, l'erreur type vaut environ $0.093$. Avec une asymétrie de $-1$ et un aplatissement de 6, elle vaut environ $0.109$, soit 16~\% de plus.

Posons $\hat{s} = \bar{X}/\hat\sigma$. Sous normalité, $\bar{X}$ et $\hat\sigma$ sont indépendants, de variances $\sigma^2/n$ et environ $\sigma^2/2n$. La méthode delta donne\index{méthode delta}
\[
\mathrm{Var}(\hat{s}) \approx \frac{1}{n} + s^2 \cdot \frac{1}{2n} = \frac{1 + s^2/2}{n}.
\]
Avec $s = 0.3$ et $n = 120$, cela donne $0.0933$.

Avec la formule donnée, le numérateur vaut $1 + 0.3 + \frac54 \cdot 0.09 = 1.41$, donc le résultat vaut $0.1085$. La formule est aussi testée par une simulation sur un mélange de rendements asymétriques et à queues épaisses, où elle est juste à 5~\% près.

Une asymétrie négative l'élargit, car elle couple la moyenne à la volatilité~: une forte perte fait à la fois baisser la moyenne et monter la volatilité, si bien que les erreurs sur les deux font varier le ratio dans le même sens. Les queues épaisses rendent l'estimation de la volatilité plus bruitée, donc le second terme est plus grand.

\iqwrong{Appliquer $1/\sqrt{n}$ sans aucune correction. Pour un ratio de Sharpe faible et des queues fines, cela reste acceptable. Sur une stratégie qui vend des options, avec une forte asymétrie négative, l'intervalle est beaucoup trop étroit.}
\iqfollow{Et en unités annuelles~?}{On multiplie par $\sqrt{12}$~: un ratio de Sharpe annuel de $1.04$ a une erreur type d'environ $0.32$ pour des rendements gaussiens, et d'environ $0.38$ dans le cas à queues épaisses. Le record se trouve alors à environ 3 erreurs types de zéro.}
\end{solution}

\begin{solution}{iq:iv1:statistics:32}
Le bon ratio de Sharpe annuel vaut environ $0.86$, et non $1.04$. Avec une autocorrélation de $-0.2$, il vaut environ $1.25$. Les stratégies peu liquides ou valorisées par modèle présentent souvent une autocorrélation positive, car leurs cours publiés évoluent lentement.

Le rendement annuel est la somme de 12 rendements mensuels, donc sa moyenne vaut $12\mu$\index{annualisation du ratio de Sharpe}, et sa variance vaut
\[
\sigma^2\Bigl(12 + 2\sum_{k = 1}^{11}(12 - k)\rho^k\Bigr).
\]
La règle usuelle ne garde que le premier terme. Avec $\rho = 0.2$, le crochet vaut environ $17.4$, et le ratio de Sharpe annuel vaut donc $0.3 \cdot 12/\sqrt{17.4} \approx 0.86$. Une simulation de 100\,000 années de ce processus donne le même résultat.

Une autocorrélation positive fait paraître la volatilité mensuelle d'une stratégie plus faible que son risque sur un long horizon, et donc un ratio de Sharpe trop élevé. Des cours périmés en sont la cause~: un actif privé valorisé par expertise, ou un portefeuille peu liquide valorisé par modèle, évolue lentement, ce qui lisse les rendements d'un mois sur l'autre. Un véritable momentum peut aussi y contribuer.

\iqwrong{Appliquer $\sqrt{12}$ dans tous les cas. Cela suppose des rendements indépendants. Si l'on voit des rendements trop lisses venant d'un portefeuille peu liquide, il faut traiter l'autocorrélation comme un signal d'alerte et corriger le lissage de la série avant de calculer le risque.}
\end{solution}

\begin{solution}{iq:iv1:statistics:33}
L'erreur type de la différence vaut environ $\sqrt{2(1 - 0.8)/5} \approx 0.283$, contre environ $0.632$ si les stratégies étaient indépendantes. La corrélation rend donc la comparaison $2.2$ fois plus précise, mais un écart de $0.28$ ne représente encore qu'une seule erreur type.

Chaque ratio de Sharpe mesuré a une variance d'environ $1/T$. La covariance de deux d'entre eux vaut environ $\rho/T$\index{comparaison de ratios de Sharpe}, car elle provient surtout de la corrélation entre les deux moyennes. La différence a donc une variance de $2(1 - \rho)/T$. Avec des données quotidiennes, le terme de volatilité est négligeable, car le ratio de Sharpe de chaque jour est petit, et une simulation sur 5 ans de données quotidiennes donne $0.283$.

Ainsi, sur 5 ans, deux stratégies aussi bonnes l'une que l'autre, avec cette corrélation, paraîtront souvent différentes~: un écart supérieur à $0.28$ en faveur de l'une d'elles survient plus de 10~\% du temps dans la simulation. Si l'on vient de choisir la meilleure de deux stratégies similaires sur la foi d'un backtest, on a probablement choisi du bruit. Le choix peut rester sans conséquence, puisqu'elles sont corrélées, mais il ne faut pas s'attendre à ce que l'écart persiste.

\iqwrong{Utiliser les deux erreurs types séparées, ou la formule à un échantillon $1/\sqrt{T} = 0.447$ pour l'écart. Les deux surestiment l'erreur de la différence, et elles laissent penser que les stratégies ne peuvent pas être distinguées alors qu'elles peuvent l'être.}
\iqwhy{Un gain constaté dans un backtest après la modification d'un paramètre est une différence entre deux stratégies fortement corrélées. Sa précision est souvent plus grande qu'on ne l'imagine, mais elle ne suffit pas à justifier la plupart des petites améliorations.}
\end{solution}

\begin{solution}{iq:iv1:statistics:34}
Le quantile empirique a une erreur type d'environ $0.118\sigma$, et $2.326\hat\sigma$ en a une d'environ $0.052\sigma$. La première est environ $2.3$ fois plus grande.

On considère le nombre de points situés sous le niveau $q$\index{erreur type d'un quantile}~: il suit une loi binomiale, d'écart type $\sqrt{np(1 - p)}$. La pente locale de son espérance vaut $nf(q)$, donc la conversion du comptage en position donne
\[
\mathrm{se}(\hat{q}) = \frac{1}{f(q)}\sqrt{\frac{p(1 - p)}{n}} = \frac{\sqrt{0.0099/1000}}{0.0267} \approx 0.118.
\]
Pour l'autre~: $\mathrm{se}(\hat\sigma) = \sigma/\sqrt{2n}$, multipliée par $2.326$. Des simulations confirment les deux résultats.

La première est plus grande parce qu'elle n'utilise que les points proches de la queue à 1~\%. Seuls une dizaine de points s'y trouvent, alors que $\hat\sigma$ les utilise tous. Dans une queue peu dense, chaque point déplace beaucoup le quantile.

Mais la seconde n'est juste que si la loi est gaussienne. Sur des rendements de loi $t$ à 3 degrés de liberté, testés également, la règle gaussienne est biaisée de plus de $0.3\sigma$, alors que le quantile empirique est presque sans biais. Il faut donc choisir le quantile empirique lorsque l'on se méfie de la forme, et en présenter le large intervalle. Une réponse précise à la mauvaise question est pire encore.

\iqwrong{Choisir une règle sur son erreur type seule. L'erreur totale inclut le biais, et en queue de distribution, l'erreur de modèle domine généralement l'erreur d'échantillonnage.}
\end{solution}

\begin{solution}{iq:iv1:statistics:35}
Pour des rendements gaussiens, l'erreur type vaut environ $0.141~\%$. Avec un aplatissement de 6, elle vaut environ $0.222~\%$, soit $\sqrt{5/2} \approx 1.58$ fois plus. Un intervalle du khi-deux a la mauvaise largeur, et il sous-couvre.

La variance de la variance empirique vaut environ $\sigma^4(\kappa - 1)/n$, où $\kappa$ est l'aplatissement\index{erreur type de la volatilité}. Une loi normale a $\kappa = 3$, donc cela vaut $2\sigma^4/n$. Pour l'écart type, la méthode delta divise par $2\sigma$~:
\[
\mathrm{se}(s) \approx \frac{\sigma}{2}\sqrt{\frac{\kappa - 1}{n}}.
\]
Pour $\kappa = 3$, cela vaut $\sigma/\sqrt{2n} = 2~\%/\sqrt{202} = 0.141~\%$. Pour $\kappa = 6$, cela vaut $1~\% \cdot \sqrt{5/101} = 0.222~\%$. Des simulations de tirages gaussiens et de Laplace, qui ont un aplatissement de 6, confirment ces valeurs.

L'intervalle du khi-deux est construit avec $\kappa = 3$. Avec des queues épaisses, la couverture réelle d'un intervalle nominal à 95~\% est, dans la simulation, inférieure à 85~\%. Il faut le corriger, par un bootstrap ou par une formule fondée sur l'aplatissement empirique, en gardant à l'esprit que l'estimation de l'aplatissement est elle-même très bruitée.

\iqwrong{Utiliser $\sigma/\sqrt{2n}$ dans tous les cas. La formule paraît universelle, car elle ne contient aucun paramètre de forme, mais c'est précisément parce qu'elle en a supposé un.}
\end{solution}

\begin{solution}{iq:iv1:statistics:36}
La vraie corrélation vaut environ $0.057$, soit près de trois fois le $0.02$ mesuré.

Supposons que $A = S + \text{bruit}$ et $B = F + \text{bruit}$, avec des bruits indépendants\index{atténuation}. Le bruit ne modifie pas la covariance, mais il augmente bien les variances. Donc
\[
\rho_{AB} = \rho_{SF}\sqrt{R_A R_B},
\]
où $R = \mathrm{Var}(\text{partie vraie})/\mathrm{Var}(\text{total})$. Ici $R_A = 0.5$ et $R_B = 0.25$, donc $\rho_{SF} = 0.02/\sqrt{0.125} \approx 0.057$. Une simulation avec 40 millions de points retrouve $0.02$ à partir de la vraie valeur $0.057$. Sans bruit, les deux sont identiques, ce qui sert de contrôle.

En finance, les rendements sont surtout du bruit, si bien qu'une corrélation brute de $0.02$ à $0.05$ peut traduire une relation forte avec la petite partie prévisible. Cela montre aussi que nettoyer le signal a de la valeur~: si l'on réduisait de moitié son bruit, la corrélation mesurée augmenterait. Cela ne signifie pas que l'on puisse trader mieux que $0.02$, puisque l'on ne peut trader que le signal bruité contre le rendement bruité.

\iqwrong{Considérer cette correction comme un moyen de gonfler sa performance mesurée. On ne peut trader que ce que l'on mesure. Le nombre corrigé renseigne sur la limite, et sur ce que des données plus propres pourraient donner.}
\end{solution}

\begin{solution}{iq:iv1:statistics:37}
La moyenne annuelle a une erreur type d'environ $6.3$ points, quelle que soit la fréquence d'échantillonnage. La volatilité annuelle en a une d'environ $0.28$ point avec des données quotidiennes, et $1.29$ avec des données mensuelles, et elle continuerait de diminuer avec un échantillonnage plus fin. On peut donc estimer la volatilité bien mieux que la moyenne, environ 22 fois mieux ici. En pratique, le bruit de microstructure limite la finesse possible de l'échantillonnage.

Pour la moyenne~: la somme des rendements sur $T$ années ne dépend que des cours de début et de fin\index{estimation de la moyenne et de la variance}, si bien qu'échantillonner plus souvent n'apporte aucune information. Son erreur type vaut $\sigma/\sqrt{T} = 20/\sqrt{10}$.

Pour la volatilité~: avec $n$ rendements, l'erreur type vaut $\sigma/\sqrt{2n}$. Pour $n = 252 \cdot 10 = 2520$, cela donne $20/\sqrt{5040} = 0.28$. Des simulations avec 120 et 2520 pas confirment les deux résultats.

Estimer des moyennes demande donc de longues histoires. Estimer la volatilité demande un échantillonnage fin.

À très fine échelle, le prix n'est pas une diffusion pure. Le rebond bid-ask et les pas de cotation discrets ajoutent du bruit, dont l'importance croît par rapport à la variance, et les rendements présentent aussi des sauts et des régularités intrajournalières. Il existe donc un intervalle d'échantillonnage en deçà duquel l'estimation se dégrade. Par ailleurs, la volatilité elle-même varie, si bien qu'un long échantillon ne répond pas à la question de ce qu'elle vaut aujourd'hui.

\iqwrong{Dire que davantage de données donne une meilleure estimation de tout. La durée et la fréquence d'échantillonnage sont deux sortes de «~plus~» différentes, et elles aident des quantités différentes.}
\end{solution}

\begin{solution}{iq:iv1:statistics:38}
Le risque vaut $d - (d - 2)^2 E[1/|X|^2]$, ce qui est inférieur à $d$ pour tout $\theta$. Pour $\theta = 0$ et $d = 10$, il vaut 2, contre 10 pour $X$.

Le lemme de Stein donne $E[(X_i - \theta_i)g(X)] = E[\partial_i g(X)]$\index{lemme de Stein}. On écrit l'estimateur sous la forme $X - h$, avec $h = (d - 2)X/|X|^2$. Alors
\[
E|X - h - \theta|^2 = d - 2\sum_i E[(X_i - \theta_i)h_i] + E|h|^2 = d - 2E[\nabla \cdot h] + E|h|^2.
\]
Un calcul direct donne $\nabla \cdot h = (d - 2)^2/|X|^2$ et $|h|^2 = (d - 2)^2/|X|^2$. Le risque vaut donc $d - 2(d - 2)^2 E[1/|X|^2] + (d - 2)^2 E[1/|X|^2]$, c'est-à-dire la formule ci-dessus. Comme $1/|X|^2 > 0$, le risque est inférieur à $d$ quel que soit $\theta$, pour $d \ge 3$. Pour $d = 2$, le facteur s'annule.

Pour $\theta = 0$, $|X|^2$ suit une loi du khi-deux, avec $E[1/|X|^2] = 1/(d - 2)$, si bien que le risque vaut $d - (d - 2) = 2$. Des simulations pour différentes valeurs de $\theta$ confirment la formule et le gain\index{estimateur de James--Stein}.

C'est étrange, car les composantes peuvent être quelconques et sans lien entre elles. Pourtant, l'erreur totale diminue si l'on rétrécit toutes les estimations vers un point commun, même si l'estimation de chacune prise isolément ne s'améliore pas. Les rendements moyens de nombreux actifs en sont un bon exemple, car ils sont bruités et proches les uns des autres. On rétrécit vers leur moyenne générale plutôt que vers 0, ce qui laisse un facteur en $d - 3$. Les optimiseurs de portefeuille obtiennent de meilleurs résultats avec de telles estimations.

\iqwrong{Dire que cela contredit l'admissibilité de la moyenne empirique en dimension un. Celle-ci tient pour $d \le 2$. Le résultat concerne la perte totale, et non une composante prise isolément.}
\end{solution}

\begin{solution}{iq:iv1:statistics:39}
L'information de Fisher par observation vaut $\tfrac12$. Ainsi, l'estimateur du maximum de vraisemblance a une variance asymptotique de $2/n$, la médiane a $\pi^2/(4n) \approx 2.47/n$, et la moyenne empirique n'a pas de variance finie et ne vaut pas mieux qu'un seul tirage, quelle que soit la valeur de $n$.

Le score vaut $\frac{2z}{1 + z^2}$, avec $z = x - \theta$\index{information de Fisher}. Donc
\[
I = \int \frac{4z^2}{(1 + z^2)^2} \cdot \frac{dz}{\pi(1 + z^2)}.
\]
En posant $z = \tan\phi$, cela devient $\frac{4}{\pi}\int_{-\pi/2}^{\pi/2} \sin^2\phi\cos^2\phi\, d\phi = \frac{4}{\pi} \cdot \frac{\pi}{8} = \tfrac12$.

La variance de la médiane vaut $1/(4nf(0)^2)$, avec $f(0) = 1/\pi$. La médiane a donc une efficacité de $8/\pi^2 \approx 0.81$ par rapport à l'estimateur du maximum de vraisemblance.

La moyenne de $n$ tirages de Cauchy suit la même loi de Cauchy qu'un seul tirage, si bien que faire la moyenne n'apporte rien. Une simulation montre que la médiane de $|\bar{X}|$ reste égale à 1.

Pour calculer l'estimateur~: la vraisemblance peut présenter plusieurs maximums locaux, donc il ne faut pas partir au hasard. On part de la médiane, qui est déjà convergente, et l'on effectue quelques pas de scoring de Fisher, en ajoutant $\frac{2}{n}\sum \frac{2z_i}{1 + z_i^2}$ à chaque pas\index{scoring de Fisher}. Cela converge vers le bon maximum pour $n$ grand. Dans une simulation avec $n = 301$, le résultat a une variance proche de $2/n$, et inférieure à celle de la médiane.

\iqwrong{Utiliser la moyenne ou un intervalle fondé sur la variance. Ni l'une ni l'autre n'existent au sens usuel pour cette loi. Une seule valeur aberrante peut déplacer une moyenne n'importe où, d'où la nécessité d'un estimateur robuste ou de la vraisemblance elle-même.}
\end{solution}

\begin{solution}{iq:iv1:statistics:40}
Soient $V$ le nombre de faux rejets et $R$ le nombre total. Pour une hypothèse nulle vraie de p-valeur $p_i$, on conditionne sur toutes les autres p-valeurs. Soit $R_i$ le nombre de rejets si $p_i$ valait 0. Alors l'hypothèse nulle $i$ est rejetée exactement lorsque $p_i \le R_i\alpha/m$, et dans ce cas le nombre total de rejets vaut exactement $R_i$.

La démonstration comporte deux étapes\index{démonstration du taux de fausses découvertes}.

Première étape~: si $p_i \le R_i\alpha/m$, alors en faisant remonter $p_i$ de 0 à sa valeur réelle, la liste triée passe encore au rang $k = R_i$, car $p_i$ est toujours sous le seuil à ce rang, et aucun $k$ plus grand ne peut passer, car il passerait aussi avec $p_i = 0$. Il y a donc $R_i$ rejets, dont l'hypothèse nulle $i$. Si $p_i > R_i\alpha/m$, alors l'hypothèse nulle $i$ n'est pas rejetée~: la rejeter supposerait un $k$ passant avec $k\alpha/m \ge p_i$, donc $k > R_i$, ce qui là encore est impossible.

Deuxième étape~: $R_i$ ne dépend que des autres p-valeurs, qui sont indépendantes de $p_i$, et $p_i$ est uniforme. Donc
\[
E\Bigl[\frac{\mathbf{1}\{i \text{ rejetée}\}}{R} \Bigm| \text{autres}\Bigr] = \frac{P(p_i \le R_i\alpha/m)}{R_i} = \frac{R_i\alpha/m}{R_i} = \frac{\alpha}{m}.
\]
Le $R_i$ se simplifie, et c'est pour cela que le résultat est si simple. On somme sur les $m_0$ hypothèses nulles vraies~:
\[
\mathrm{FDR} = E\Bigl[\frac{V}{\max(R, 1)}\Bigr] = \sum_{i \text{ nulle}} E\Bigl[\frac{\mathbf{1}\{i \text{ rejetée}\}}{R}\Bigr] = \frac{m_0\alpha}{m}.
\]
C'est exactement cette valeur, et non seulement une borne, sous indépendance. Une simulation avec $m = 20$, $m_0 = 15$ et $\alpha = 0.1$ donne $0.075$, comme annoncé. Un test direct de la première étape, qui fait varier une p-valeur sur une grille tandis que les autres restent fixes, donne une probabilité de rejet de $R_i\alpha/m$.

\iqwrong{Chercher à conditionner sur le nombre de rejets $R$. Celui-ci dépend de $p_i$, donc la condition s'embrouille. L'astuce consiste à utiliser $R_i$, qui ne dépend pas de $p_i$.}
\end{solution}
